% Generated draft: requires mathematical and source review.
\documentclass[11pt]{article}
\usepackage[margin=25mm]{geometry}
\usepackage{fontspec}
\setmainfont{Latin Modern Roman}
\usepackage{amsmath,amssymb,mathtools,unicode-math}
\setmathfont{Latin Modern Math}
\usepackage{enumitem,needspace,xurl,hyperref}
\setlength{\parindent}{0pt}
\setlength{\parskip}{5pt}
\setlength{\emergencystretch}{3em}
\newcommand{\N}{\mathbb{N}}
\newcommand{\Z}{\mathbb{Z}}
\newcommand{\Q}{\mathbb{Q}}
\newcommand{\R}{\mathbb{R}}
\newcommand{\C}{\mathbb{C}}
\newcommand{\F}{\mathbb{F}}
\newcommand{\A}{\mathbb{A}}
\newcommand{\PP}{\mathbb{P}}
\newcommand{\E}{\mathbb{E}}
\newcommand{\eps}{\varepsilon}
\newcommand{\dd}{\,\mathrm{d}}
\begin{document}
\section*{UnsolvedMath ID: 30000698}
\textbf{Source problem code:} catalogue entry
\subsection{Definitions and mathematical statement}
Let $f:\C\to\C$ be a transcendental entire function, meaning an entire function that is not a polynomial, and let $N_f$ be the meromorphic Newton map given by $N_f(z)=z-f(z)/f'(z)$ wherever the quotient is defined, with removable singularities filled in. Suppose that $N_f$ has an invariant Baker domain $U$. In the standard terminology, $U$ is a component of the Fatou set, $N_f(U)\subseteq U$, and $N_f^{\circ m}\to\infty$ locally uniformly on $U$ as $m\to\infty$. The supplied statement abbreviates this by invariance and escape of the iterates. Put $M_f(r)=\max_{|z|=r}|f(z)|$ and define the order, allowing the value $+\infty$, by \[\rho(f)=\limsup_{r\to\infty}\frac{\log\log M_f(r)}{\log r}.\] The source asks the following two questions for every such $f$. First, must there exist $K>0$ and a continuous curve $\gamma:[0,\infty)\to\C$ with $|\gamma(t)|\to\infty$ such that \[|f(\gamma(t))|\leq |\gamma(t)|^K\qquad(t\geq0)?\] Equivalently, one may ask for this inequality on a tail of a curve and then discard its initial segment. No requirement that $\gamma$ lie in $U$ is stated. Second, must $\rho(f)\geq\tfrac12$? These are separate questions. The supplied literature triage, dated 2026-08-21, labels the record partially solved. It reports that a related Douady question about asymptotic values is false in general, that the threshold $1/2$ is sharp for known Baker-domain constructions, and that the universal polynomial-growth-on-a-curve question remains unsettled in the sources located. Those reports are not an independently verified answer to either displayed question; in particular, sharpness for constructions alone does not specify a proof of the universal lower bound.
\subsection{Short English statement}
For a transcendental entire function whose Newton map has an invariant Baker domain, ask whether the function grows at most polynomially along some curve to infinity, and whether its order is at least $1/2$. Preserve the supplied 2026-08-21 partial-status assessment and its distinction from the related asymptotic-value question.
\subsection{Sources}
\begin{itemize}\item Local UnsolvedMath registry, permanent ID 30000698, source problem catalogue entry, \emph{Growth and Order of Entire Functions with Newton Baker Domains}, proposed year 2007; no individual proposer supplied.\item \emph{Normal Families and Complex Dynamics} (2007), Oberwolfach Reports, DOI 10.4171/owr/2007/09, pp.\ 487--548: \url{https://doi.org/10.4171/owr/2007/09}. Supplied source and literature label: original problem report.\item Supplied literature label: Baker domains for Newton method, \url{https://doi.org/10.5802/aif.2277}.\item Supplied literature label: counterexample without finite asymptotic values, \url{https://arxiv.org/abs/math/0601496}.\end{itemize}
\end{document}
